\chapter{1992年全国卷（文）}
\section{单选题}

\begin{problem}
\(\frac{\log_8 9}{\log_2 3}\) 的值是（\quad）

\choices
    {\(\frac{2}{3} \)}
    {\(1\)}
    {\(\frac{3}{2} \)}
    {\(2\)}

\begin{answer}
A
\end{answer}

\begin{solution}
由换底公式，\(\log_8 9=\frac{2\log 3}{3\log 2}\)，且\(\log_2 3=\frac{\log 3}{\log 2}\)，因此
\(\frac{\log_8 9}{\log_2 3}=\frac{2\log 3}{3\log 2}\cdot\frac{\log 2}{\log 3}=\frac{2}{3}\). 故选 A.
\end{solution}
\end{problem}

\begin{problem}
已知椭圆 \(\frac{x^2}{25} + \frac{y^2}{16} = 1\) 上一点 \(P\) 到椭圆一个焦点的距离为 \(3\)，则 \(P\) 到另一焦点的距离为（\quad）

\choices
    {\(2\)}
    {\(3\)}
    {\(5\)}
    {\(7\)}

\begin{answer}
D
\end{answer}

\begin{solution}
椭圆的长半轴为 \(a=5\)，所以椭圆上任一点到两个焦点的距离之和为 \(2a=10\). 已知其中一个距离为 \(3\)，故到另一个焦点的距离为 \(10-3=7\). 故选 D.
\end{solution}
\end{problem}

\begin{problem}
如果函数 \(y=\sin(\omega x)\cos(\omega x)\) 的最小正周期是 \(4\pi\)，那么常数 \(\omega\) 为（\quad）

\choices
    {\(4\)}
    {\(2\)}
    {\(\frac{1}{2} \)}
    {\(\frac{1}{4} \)}

\begin{answer}
D
\end{answer}

\begin{solution}
利用二倍角公式，\(y=\frac{1}{2}\sin(2\omega x)\). 因而其最小正周期为 \(\frac{2\pi}{2|\omega|}=\frac{\pi}{|\omega|}\). 由题设 \(\frac{\pi}{|\omega|}=4\pi\)，得 \(|\omega|=\frac{1}{4}\). 按三角函数中角频率取正值的约定，\(\omega=\frac{1}{4}\)，故选 D.
\end{solution}
\end{problem}

\begin{problem}
在 \(\left(\frac{x}{2}-\frac{1}{\sqrt[3]{x}}\right)^8\) 的展开式中常数项是（\quad）

\choices
    {\(-28\)}
    {\(-7\)}
    {\(7\)}
    {\(28\)}

\begin{answer}
C
\end{answer}

\begin{solution}
取展开式中第二项因子的次数为 \(k\)，其一般项为 \(\mathrm C_8^k\left(\frac{x}{2}\right)^{8-k}\left(-x^{-\frac{1}{3}}\right)^k\)，其中 \(x\) 的指数为 \(8-k-\frac{k}{3}=8-\frac{4k}{3}\). 常数项要求 \(8-\frac{4k}{3}=0\)，所以 \(k=6\). 此时常数项为 \(\mathrm C_8^6\left(\frac{1}{2}\right)^2(-1)^6=28\cdot\frac{1}{4}=7\). 故选 C.
\end{solution}
\end{problem}

\begin{problem}
已知轴截面是正方形的圆柱的高与球的直径相等，则圆柱的全面积与球的表面积的比是（\quad）

\choices
    {\(6:5\)}
    {\(5:4\)}
    {\(4:3\)}
    {\(3:2\)}

\begin{answer}
D
\end{answer}

\begin{solution}
设圆柱底面半径为 \(r\)，高为 \(h\). 轴截面为正方形，所以 \(h=2r\). 又圆柱高等于球的直径，故球半径也是 \(r\). 圆柱全面积为 \(2\pi r^2+2\pi rh=2\pi r^2+4\pi r^2=6\pi r^2\)，球表面积为 \(4\pi r^2\)，所以所求比为 \(6\pi r^2:4\pi r^2=3:2\). 故选 D.
\end{solution}
\end{problem}

\begin{problem}
如图，图中曲线是幂函数 \(y=x^n\) 在第一象限的图象. 已知 \(n\) 取 \(\pm 2\)，\(\pm \frac{1}{2}\) 四个值，则相应于曲线 \(C_{1}\)，\(C_{2}\)，\(C_{3}\)，\(C_{4}\) 的 \(n\) 依次为（\quad）

\bitmapfigure[max width=0.28\linewidth]{img/1992/national_liberal/q6_fig1.png}

\choices
    {\(-2,-\frac{1}{2},\frac{1}{2},2\)}
    {\(2\)，\(\frac{1}{2}\)，\(-\frac{1}{2}\)，\(-2\)}
    {\(-\frac{1}{2}\)，\(-2\)，\(2\)，\(\frac{1}{2}\)}
    {\(2\)，\(\frac{1}{2}\)，\(-2\)，\(-\frac{1}{2}\)}

\begin{answer}
B
\end{answer}

\begin{solution}
当 \(x>1\) 时，正指数幂函数递增，且 \(x^2>x^{1/2}\)，所以图中上面的递增曲线 \(C_1\) 对应 \(n=2\)，下面的递增曲线 \(C_2\) 对应 \(n=\frac{1}{2}\). 负指数幂函数递减，且 \(x^{-1/2}>x^{-2}\)，所以较上面的递减曲线 \(C_3\) 对应 \(n=-\frac{1}{2}\)，最下面的递减曲线 \(C_4\) 对应 \(n=-2\). 故依次为 \(2\)，\(\frac{1}{2}\)，\(-\frac{1}{2}\)，\(-2\)，选 B.
\end{solution}
\end{problem}

\begin{problem}
若 \(\log_a 2 < \log_b 2 < 0\)，则（\quad）

\choices
    {\(0 < a < b < 1\)}
    {\(0 < b < a < 1\)}
    {\(a > b > 1\)}
    {\(b > a > 1\)}

\begin{answer}
B
\end{answer}

\begin{solution}
因为 \(\log_a2<0\) 且 \(\log_b2<0\)，而真数 \(2>1\)，所以 \(0<a<1\)，\(0<b<1\). 在 \(0<t<1\) 上，\(g(t)=\log_t2=\frac{\ln2}{\ln t}\) 是减函数，因为 \(g'(t)=-\frac{\ln2}{t(\ln t)^2}<0\). 由 \(g(a)<g(b)\) 得 \(a>b\)，所以 \(0<b<a<1\). 故选 B.
\end{solution}
\end{problem}

\begin{problem}
原点关于直线 \(8x + 6y = 25\) 的对称点坐标为（\quad）

\choices
    {\(\left(2,\frac{3}{2}\right)\)}
    {\(\left(\frac{25}{8}, \frac{25}{6}\right)\)}
    {\(\left(3,4\right)\)}
    {\(\left(4,3\right)\)}

\begin{answer}
D
\end{answer}

\begin{solution}
原点到直线的垂足在法向量 \((8,6)\) 方向上，设垂足为 \((8t,6t)\). 代入直线方程得 \(8\cdot8t+6\cdot6t=25\)，即 \(100t=25\)，所以垂足为 \(\left(2,\frac{3}{2}\right)\). 对称点是原点关于垂足的中心对称点，故坐标为 \(\left(4,3\right)\). 故选 D.
\end{solution}
\end{problem}

\begin{problem}
在四棱锥的四个侧面中，直角三角形最多可有（\quad）

\choices
    {\(1\) 个}
    {\(2\) 个}
    {\(3\) 个}
    {\(4\) 个}

\begin{answer}
D
\end{answer}

\begin{solution}
四棱锥只有四个侧面，所以直角三角形的个数不超过 \(4\). 下面说明四个可以同时成立：取正方形 \(ABCD\) 为底面，取顶点 \(S\) 在底面顶点 \(A\) 的垂线上. 则在侧面 \(SAB\)、\(SAD\) 中，\(SA\perp AB\)、\(SA\perp AD\)，它们在 \(A\) 处为直角；在侧面 \(SBC\) 中，\(SB\) 的水平投影平行于 \(AB\)，而 \(BC\perp AB\)，故 \(SB\perp BC\)；同理在侧面 \(SCD\) 中 \(SD\perp DC\). 因而四个侧面都可以是直角三角形，最大值为 \(4\). 故选 D.
\end{solution}
\end{problem}

\begin{problem}
圆心在抛物线 \(y^{2} = 2x \) 上，且与 \(x \) 轴和该抛物线的准线都相切的一个圆的方程是（\quad）

\choices
    {\(x^{2} + y^{2} - x - 2y - \frac{1}{4} = 0\)}
    {\(x^{2} + y^{2} + x - 2y + 1 = 0\)}
    {\(x^{2} + y^{2} - x - 2y + 1 = 0\)}
    {\(x^{2} + y^{2} - x - 2y + \frac{1}{4} = 0\)}

\begin{answer}
D
\end{answer}

\begin{solution}
将抛物线写成 \(y^2=4px\)，得 \(p=\frac{1}{2}\)，所以准线为 \(x=-\frac{1}{2}\). 设圆心为 \((u,v)\)，半径为 \(r\). 圆同时与 \(x\) 轴和准线相切，因此 \(r=|v|=\left|u+\frac{1}{2}\right|\). 又圆心在抛物线上，所以 \(v^2=2u\). 消去半径得
\(2u=\left(u+\frac{1}{2}\right)^2\)，即 \(\left(u-\frac{1}{2}\right)^2=0\)，从而 \(u=\frac{1}{2}\)，\(v=\pm1\). 题目选项列出的是上半平面圆，取圆心 \(\left(\frac{1}{2},1\right)\)，半径为 \(1\)，其方程为 \(\left(x-\frac{1}{2}\right)^2+(y-1)^2=1\)，即 \(x^2+y^2-x-2y+\frac{1}{4}=0\). 故选 D.
\end{solution}
\end{problem}

\begin{problem}
在 \([0,2\pi]\) 上满足 \(\sin x \geqslant \frac{1}{2}\) 的 \(x\) 的取值范围是（\quad）

\choices
    {\(\left[0, \frac{\pi}{6}\right]\)}
    {\(\left[\frac{\pi}{6}, \frac{5\pi}{6}\right]\)}
    {\(\left[\frac{\pi}{6}, \frac{2\pi}{3}\right]\)}
    {\(\left[\frac{5\pi}{6}, \pi\right]\)}

\begin{answer}
B
\end{answer}

\begin{solution}
在一个周期内，\(\sin x=\frac{1}{2}\) 的两个交点为 \(x=\frac{\pi}{6}\) 和 \(x=\frac{5\pi}{6}\). 正弦函数在这两个点之间不小于 \(\frac{1}{2}\)，而在其余部分小于该值，所以在 \([0,2\pi]\) 上的解集为 \(\left[\frac{\pi}{6},\frac{5\pi}{6}\right]\). 故选 B.
\end{solution}
\end{problem}

\begin{problem}
已知直线 \(l_{1}\) 和 \(l_{2}\) 夹角的平分线为 \(y = x\)，如果 \(l_{1}\) 的方程是 \(ax + by + c = 0\)（\(ab > 0\)），那么 \(l_{2}\) 的方程是（\quad）

\choices
    {\(bx + ay + c = 0\)}
    {\(ax - by + c = 0\)}
    {\(bx + ay - c = 0\)}
    {\(bx - ay + c = 0\)}

\begin{answer}
A
\end{answer}

\begin{solution}
关于直线 \(y=x\) 对称时，点的坐标互换，因此直线 \(l_1\colon ax+by+c=0\) 的对称直线为 \(ay+bx+c=0\)，即 \(bx+ay+c=0\). 这两条直线关于 \(y=x\) 对称，故 \(y=x\) 是它们的一个夹角平分线，选项 A 正确.
\end{solution}
\end{problem}

\begin{problem}
如果 \(\alpha\)，\(\beta \in \left(\frac{\pi}{2}, \pi\right)\) 且 \(\tan \alpha < \cot \beta\)，那么必有（\quad）

\choices
    {\(\alpha < \beta\)}
    {\(\beta < \alpha\)}
    {\(\alpha + \beta < \frac{3}{2}\pi\)}
    {\(\alpha + \beta > \frac{3}{2}\pi\)}

\begin{answer}
C
\end{answer}

\begin{solution}
令 \(u=\alpha-\pi\)，\(v=\frac{\pi}{2}-\beta\)，则 \(u\)，\(v\in\left(-\frac{\pi}{2},0\right)\)，并且 \(\tan u=\tan\alpha\)，\(\tan v=\cot\beta\). 正切函数在 \(\left(-\frac{\pi}{2},\frac{\pi}{2}\right)\) 上递增，所以由 \(\tan u<\tan v\) 得 \(u<v\). 因而 \(\alpha-\pi<\frac{\pi}{2}-\beta\)，即 \(\alpha+\beta<\frac{3\pi}{2}\). 故选 C.
\end{solution}
\end{problem}

\begin{problem}
如图，在棱长为 \(1\) 的正方体 \(ABCD-A_{1}B_{1}C_{1}D_{1}\) 中，\(M\) 和 \(N\) 分别为 \(A_{1}B_{1}\) 和 \(BB_{1}\) 的中点，那么直线 \(AM\) 与 \(CN\) 所成角的余弦值是（\quad）

\bitmapfigure[max width=5.3cm]{img/1992/national_liberal/q14_fig1.png}

\choices
    {\(\frac{\sqrt{3}}{2}\)}
    {\(\frac{\sqrt{10}}{10}\)}
    {\(\frac{3}{5}\)}
    {\(\frac{2}{5}\)}

\begin{answer}
D
\end{answer}

\begin{solution}
建立直角坐标系，取 \(A=(0,0,0)\)，\(B=(1,0,0)\)，\(C=(1,1,0)\)，\(A_1=(0,0,1)\)，\(B_1=(1,0,1)\). 则 \(M=\left(\frac{1}{2},0,1\right)\)，\(N=\left(1,0,\frac{1}{2}\right)\). 取两直线方向向量 \(\overrightarrow{AM}=\left(\frac{1}{2},0,1\right)\) 和 \(\overrightarrow{CN}=\left(0,-1,\frac{1}{2}\right)\)，有
\(\overrightarrow{AM}\cdot\overrightarrow{CN}=\frac{1}{2}\)，\(\left|\overrightarrow{AM}\right|=\frac{\sqrt{5}}{2}\)，\(\left|\overrightarrow{CN}\right|=\frac{\sqrt{5}}{2}\). 所成角取锐角，故余弦值为 \(\frac{\frac{1}{2}}{\frac{\sqrt{5}}{2}\cdot\frac{\sqrt{5}}{2}}=\frac{2}{5}\). 故选 D.
\end{solution}
\end{problem}

\begin{problem}
已知复数 \(z\) 的模为 \(2\)，则 \(|z - \i|\) 的最大值为（\quad）

\choices
    {\(1\)}
    {\(2\)}
    {\(\sqrt{5}\)}
    {\(3\)}

\begin{answer}
D
\end{answer}

\begin{solution}
把复数看作复平面上的点，\(|z|=2\)，而 \(|\i|=1\). 由三角不等式，\(|z-\i|\leqslant |z|+|\i|=3\). 当 \(z=-2\i\) 时，\(|z|=2\)，且 \(|z-\i|=|-3\i|=3\)，等号可以取得. 因此最大值为 \(3\)，故选 D.
\end{solution}
\end{problem}

\begin{problem}
函数 \(y=\frac{\e^x-\e^{-x}}{2}\) 的反函数（\quad）

\choices
    {是奇函数，它在 \((0, +\infty)\) 上是减函数}
    {是偶函数，它在 \((0, +\infty)\) 上是减函数}
    {是奇函数，它在 \((0, +\infty)\) 上是增函数}
    {是偶函数，它在 \((0, +\infty)\) 上是增函数}

\begin{answer}
C
\end{answer}

\begin{solution}
原函数为 \(f(x)=\frac{\e^x-\e^{-x}}{2}\)，其定义域为 \(\R\). 有 \(f(-x)=-f(x)\)，所以它是奇函数；又 \(f'(x)=\frac{\e^x+\e^{-x}}{2}>0\)，所以它在 \(\R\) 上严格递增. 此外，\(\displaystyle\lim\limits_{x\to+\infty}f(x)=+\infty\)，\(\displaystyle\lim\limits_{x\to-\infty}f(x)=-\infty\)，故 \(f\) 的值域为 \(\R\)，反函数 \(g=f^{-1}\) 存在且定义域为 \(\R\). 若 \(y=f(x)\)，则 \(-y=f(-x)\)，从而 \(g(-y)=-x=-g(y)\)，所以 \(g\) 是奇函数. 由于 \(f\) 严格递增，\(g\) 也严格递增，因而 \(g\) 在 \((0,+\infty)\) 上是增函数，故选 C.
\end{solution}
\end{problem}

\begin{problem}
如果函数 \(f(x)=x^2+bx+c\) 对任意实数 \(t\) 都有 \(f(2+t)=f(2-t)\)，那么（\quad）

\choices
    {\(f(2) < f(1) < f(4)\)}
    {\(f(1) < f(2) < f(4)\)}
    {\(f(2) < f(4) < f(1)\)}
    {\(f(4) < f(2) < f(1)\)}

\begin{answer}
A
\end{answer}

\begin{solution}
条件说明二次函数的图象关于直线 \(x=2\) 对称，所以其对称轴为 \(x=-\frac{b}{2}=2\)，从而 \(b=-4\). 因此 \(f(x)=x^2-4x+c=(x-2)^2+c-4\). 于是 \(f(2)=c-4\)，\(f(1)=c-3\)，\(f(4)=c\)，故 \(f(2)<f(1)<f(4)\). 故选 A.
\end{solution}
\end{problem}

\begin{problem}
已知长方体的全面积为 \(11\)，\(12\) 条棱长度之和为 \(24\)，则这个长方体的一条对角线长为（\quad）

\choices
    {\(2\sqrt{3}\)}
    {\(\sqrt{14}\)}
    {\(5\)}
    {\(6\)}

\begin{answer}
C
\end{answer}

\begin{solution}
设长方体的三条棱长为 \(a\)，\(b\)，\(c>0\). 由全面积和棱长总和分别得 \(2(ab+bc+ca)=11\)，\(4(a+b+c)=24\)，即 \(ab+bc+ca=\frac{11}{2}\)，\(a+b+c=6\). 对角线长 \(d\) 满足
\(d^2=a^2+b^2+c^2=(a+b+c)^2-2(ab+bc+ca)=36-11=25\). 因为 \(d>0\)，所以 \(d=5\). 故选 C.
\end{solution}
\end{problem}

\section{填空题}

\begin{problem}
\(\displaystyle\lim\limits_{n\to \infty}\left[\frac{1}{3}-\frac{1}{9}+\frac{1}{27}+\cdots+(-1)^{n-1}\frac{1}{3^n}\right]\) 的值为\fillinblank{}.

\begin{answer}
\(\frac{1}{4}\)
\end{answer}

\begin{solution}
括号内是首项为 \(\frac{1}{3}\)、公比为 \(-\frac{1}{3}\) 的等比数列前 \(n\) 项和，因此
\(S_n=\frac{\frac{1}{3}\left(1-\left(-\frac{1}{3}\right)^n\right)}{1-\left(-\frac{1}{3}\right)}=\frac{1}{4}\left(1-\left(-\frac{1}{3}\right)^n\right)\). 因为 \(\left(-\frac{1}{3}\right)^n\to0\)，所以原极限为 \(\frac{1}{4}\).
\end{solution}
\end{problem}

\begin{problem}
已知 \(\alpha\) 在第三象限且 \(\tan \alpha = 2\)，则 \(\cos \alpha\) 的值是\fillinblank{}.

\begin{answer}
\(-\frac{\sqrt{5}}{5}\)
\end{answer}

\begin{solution}
由恒等式 \(1+\tan^2\alpha=\sec^2\alpha\)，得 \(\sec^2\alpha=1+2^2=5\)，所以 \(|\cos\alpha|=\frac{1}{\sqrt{5}}\). 因为 \(\alpha\) 在第三象限，\(\cos\alpha<0\)，故 \(\cos\alpha=-\frac{1}{\sqrt{5}}=-\frac{\sqrt{5}}{5}\).
\end{solution}
\end{problem}

\begin{problem}
方程 \(\frac{1 + 3^{-x}}{1 + 3^x} = 3\) 的解是\fillinblank{}.

\begin{answer}
\(x=-1\)
\end{answer}

\begin{solution}
令 \(t=3^x>0\)，则 \(3^{-x}=\frac{1}{t}\). 原方程化为 \(\frac{1+\frac{1}{t}}{1+t}=3\). 由于 \(t>0\)，分母不为零，整理得 \(1+\frac{1}{t}=3+3t\)，即 \(3t^2+2t-1=0\). 因式分解得 \((3t-1)(t+1)=0\). 结合 \(t>0\)，得 \(t=\frac{1}{3}\)，所以 \(3^x=3^{-1}\)，解得 \(x=-1\).
\end{solution}
\end{problem}

\begin{problem}
设含有 \(10\) 个元素的集合的全部子集数为 \(S\)，其中由 \(3\) 个元素组成的子集数为 \(T\)，则 \(\frac{T}{S}\) 的值为\fillinblank{}.

\begin{answer}
\(\frac{15}{128}\)
\end{answer}

\begin{solution}
含有 \(10\) 个元素的集合的每个元素都有“取入子集”或“不取入子集”两种选择，所以 \(S=2^{10}=1024\). 从其中选出 \(3\) 个元素组成子集的方法数为 \(T=\mathrm C_{10}^{3}=\frac{10\cdot9\cdot8}{3\cdot2\cdot1}=120\). 因而 \(\frac{T}{S}=\frac{120}{1024}=\frac{15}{128}\).
\end{solution}
\end{problem}

\begin{problem}
焦点为 \(F_{1}(-2,0)\) 和 \(F_{2}(6,0)\)，离心率为 \(2\) 的双曲线的方程是\fillinblank{}.

\begin{answer}
\(\frac{\left(x-2\right)^2}{4}-\frac{y^2}{12}=1\)
\end{answer}

\begin{solution}
两焦点中点为 \((2,0)\)，焦半距为 \(c=4\). 由离心率 \(e=\frac{c}{a}=2\)，得半实轴 \(a=2\). 双曲线满足 \(c^2=a^2+b^2\)，所以 \(b^2=4^2-2^2=12\). 因焦点在水平轴上，方程为 \(\frac{(x-2)^2}{4}-\frac{y^2}{12}=1\).
\end{solution}
\end{problem}

\section{解答题}

\begin{problem}
求 \(\sin^2 20^\circ + \cos^2 80^\circ + \sqrt{3} \sin 20^\circ \cos 80^\circ\) 的值.

\begin{solution}
利用降幂公式和积化和差公式，原式为 \(\frac{1-\cos40^\circ}{2}+\frac{1+\cos160^\circ}{2}+\frac{\sqrt3}{2}\left(\sin100^\circ+\sin(-60^\circ)\right)\).
由于 \(\cos160^\circ=-\cos20^\circ\)，\(\sin100^\circ=\sin80^\circ\)，\(\sin(-60^\circ)=-\frac{\sqrt3}{2}\)，所以原式
为 \(\frac{1}{4}-\frac{\cos20^\circ+\cos40^\circ}{2}+\frac{\sqrt3}{2}\sin80^\circ\).
又 \(\cos20^\circ+\cos40^\circ=2\cos30^\circ\cos10^\circ=\sqrt3\cos10^\circ=\sqrt3\sin80^\circ\)，故原式为 \(\frac{1}{4}\).
\end{solution}
\end{problem}

\begin{problem}
已知 \(z \in \mathbf{C}\)，解方程：\(z - 2|z| = -7 + 4\i\).

\begin{solution}
设 \(z=x+y\i\)，并令 \(r=|z|=\sqrt{x^2+y^2}\geqslant0\). 比较实部和虚部得 \(x-2r=-7\)，\(y=4\)，即 \(x=2r-7\). 代入模的关系式，得
\(r^2=(2r-7)^2+4^2=4r^2-28r+65\)，即 \(3r^2-28r+65=0\). 其判别式为 \(\Delta=(-28)^2-4\cdot3\cdot65=4\)，所以
\(r=\frac{28\pm\sqrt{4}}{2\cdot3}=5\) 或 \(r=\frac{13}{3}\). 两个根都满足 \(r\geqslant0\). 当 \(r=5\) 时，\(x=2\cdot5-7=3\)，且 \(\sqrt{x^2+4^2}=5=r\)，得 \(z=3+4\i\)；当 \(r=\frac{13}{3}\) 时，\(x=2\cdot\frac{13}{3}-7=\frac{5}{3}\)，且 \(\sqrt{\left(\frac{5}{3}\right)^2+4^2}=\frac{13}{3}=r\)，得 \(z=\frac{5}{3}+4\i\). 两个值均满足 \(r=|z|\) 及实、虚部方程，因而均为原方程的解，故 \(z=3+4\i\) 或 \(z=\frac{5}{3}+4\i\).
\end{solution}
\end{problem}

\begin{problem}
如图，已知 \(ABCD-A_{1}B_{1}C_{1}D_{1}\) 是棱长为 \(a\) 的正方体，\(E\)，\(F\) 分别为棱 \(AA_{1}\) 与 \(CC_{1}\) 的中点，求四棱锥 \(A_{1}-EBFD_{1}\) 的体积.

\bitmapfigure[max width=0.28\linewidth]{img/1992/national_liberal/q26_fig1.png}

\begin{solution}
取坐标 \(A=(0,0,0)\)，\(B=(a,0,0)\)，\(C=(a,a,0)\)，\(D=(0,a,0)\)，\(A_1=(0,0,a)\). 则 \(E=\left(0,0,\frac{a}{2}\right)\)，\(F=\left(a,a,\frac{a}{2}\right)\)，\(D_1=(0,a,a)\). 有
\(\overrightarrow{EB}=a\left(1,0,-\frac12\right)\)，\(\overrightarrow{EF}=a(1,1,0)\)，\(\overrightarrow{ED_1}=a\left(0,1,\frac12\right)\). 计算得
\(\overrightarrow{EB}\times\overrightarrow{EF}=a^2\left(\frac12,-\frac12,1\right)\)，从而 \(\left(\overrightarrow{EB}\times\overrightarrow{EF}\right)\cdot\overrightarrow{ED_1}=0\)，所以 \(E\)，\(B\)，\(F\)，\(D_1\) 共面. 又 \(\overrightarrow{EB}+\overrightarrow{FD_1}=0\)，\(\overrightarrow{BF}+\overrightarrow{D_1E}=0\)，故 \(EBFD_1\) 是平行四边形，其对角线 \(EF\) 将底面分成三角形 \(EBF\) 和 \(EFD_1\).

以 \(E\) 为起点，\(\overrightarrow{EA_1}=a\left(0,0,\frac12\right)\). 因而
\(\left|\left(\overrightarrow{EB}\times\overrightarrow{EF}\right)\cdot\overrightarrow{EA_1}\right|=\frac{a^3}{2}\)，所以三棱锥 \(A_1-EBF\) 的体积为 \(\frac{1}{6}\cdot\frac{a^3}{2}=\frac{a^3}{12}\). 另外，
\(\overrightarrow{EF}\times\overrightarrow{ED_1}=a^2\left(\frac12,-\frac12,1\right)\)，故
\(\left|\left(\overrightarrow{EF}\times\overrightarrow{ED_1}\right)\cdot\overrightarrow{EA_1}\right|=\frac{a^3}{2}\)，三棱锥 \(A_1-EFD_1\) 的体积也为 \(\frac{a^3}{12}\). 因此四棱锥体积为 \(\frac{a^3}{12}+\frac{a^3}{12}=\frac{a^3}{6}\).
\end{solution}
\end{problem}

\begin{problem}
在 \(\triangle ABC\) 中，\(BC\) 边上的高所在直线的方程为 \(x-2y+1=0\)，\(\angle A\) 的平分线所在直线的方程为 \(y=0\)，若点 \(B\) 的坐标为 \(\left(1,2\right)\)，求点 \(A\) 和点 \(C\) 的坐标.

\begin{solution}
点 \(A\) 在 \(\angle A\) 的平分线 \(y=0\) 上，也在从 \(A\) 向 \(BC\) 作出的高所在直线 \(x-2y+1=0\) 上. 联立得 \(y=0\)，\(x+1=0\)，所以 \(A=(-1,0)\).

直线 \(x-2y+1=0\) 的斜率为 \(\frac{1}{2}\)，因此 \(BC\) 的斜率为 \(-2\). 由于 \(B=(1,2)\)，得 \(BC\) 的方程为 \(y-2=-2(x-1)\)，即 \(y=-2x+4\). 直线 \(y=0\) 是 \(\angle A\) 的平分线，所以射线 \(AC\) 是射线 \(AB\) 关于 \(x\) 轴的对称射线. 直线 \(AB\) 的方程为 \(y=x+1\)，关于 \(x\) 轴对称后，直线 \(AC\) 的方程为 \(y=-x-1\). 联立 \(y=-2x+4\) 与 \(y=-x-1\)，得 \(x=5\)，\(y=-6\). 因而 \(C=(5,-6)\)，所求为 \(A=(-1,0)\)，\(C=(5,-6)\).
\end{solution}
\end{problem}

\begin{problem}
设等差数列 \(\{a_{n}\}\) 的前 \(n\) 项和为 \(S_{n}\). 已知 \(a_{3}=12\)，\(S_{12}>0\)，\(S_{13}<0\).

\begin{enumerate}
\item 求公差 \(d\) 的取值范围；
\item 指出 \(S_{1}\)，\(S_{2}\)，\(\cdots\)，\(S_{12}\) 中哪一个值最大，并说明理由.
\end{enumerate}
\begin{solution}
\begin{enumerate}
\item 由 \(a_3=a_1+2d=12\)，得 \(a_1=12-2d\). 因而
\(S_{12}=\frac{12}{2}(a_1+a_{12})=6\left(24+7d\right)>0\)，得 \(d>-\frac{24}{7}\)；又
\(S_{13}=\frac{13}{2}(a_1+a_{13})=\frac{13}{2}\left(24+8d\right)<0\)，得 \(d<-3\). 所以 \(-\frac{24}{7}<d<-3\).
\item 等差数列的项为 \(a_n=a_3+(n-3)d=12+(n-3)d\). 由 \(d>-\frac{24}{7}>-4\)，得 \(a_6=12+3d>12-12=0\)；由 \(d<-3\)，得 \(a_7=12+4d<12-12=0\). 又因 \(d<0\)，数列各项递减，所以 \(a_1\)，\(\ldots\)，\(a_6>0\)，而 \(a_7\)，\(\ldots\)，\(a_{12}<0\). 因为 \(S_{n+1}-S_n=a_{n+1}\)，故 \(S_1<S_2<\cdots<S_6\)，且 \(S_6>S_7>\cdots>S_{12}\). 因此 \(S_6\) 最大.
\end{enumerate}
\end{solution}
\end{problem}
